\section{Results}
\label{sec:results}
% =================================================================

\subsection{Platform Architecture Overview}

The PAI Oncology Trial FL platform follows a layered architecture where five infrastructure pillars support the end-to-end clinical trial workflow. Table~\ref{tab:architecture} maps each pillar to its codebase location, version of introduction, and quantitative scope.

\begin{table}[H]
\centering
\caption{Five-pillar architecture with codebase metrics}
\label{tab:architecture}
\small
\begin{tabular}{@{}clcrl@{}}
\toprule
\textbf{\#} & \textbf{Pillar} & \textbf{Directory} & \textbf{LOC} & \textbf{Version} \\
\midrule
1 & Privacy Infrastructure & \texttt{privacy/} & $\sim$5,400 & v0.6.0 \\
2 & Regulatory Infrastructure & \texttt{regulatory/}, \texttt{reg-sub/} & $\sim$10,000 & v0.6.0, v0.9.1 \\
3 & Cross-Framework Unification & \texttt{unification/} & $\sim$8,200 & v0.3.0 \\
4 & Standards \& Benchmarking & \texttt{q1-2026-standards/} & $\sim$3,800 & v0.3.0 \\
5 & Multi-Org Cooperation & \texttt{federated/} & $\sim$12,000 & v0.1.0--v0.2.0 \\
\midrule
  & Supporting: Clinical Analytics & \texttt{clinical-analytics/} & $\sim$4,900 & v0.9.0 \\
  & Supporting: Digital Twins & \texttt{digital-twins/} & $\sim$4,500 & v0.4.0 \\
  & Supporting: Physical AI & \texttt{physical\_ai/} & $\sim$6,200 & v0.1.0 \\
  & Supporting: Agentic AI & \texttt{agentic-ai/} & $\sim$6,700 & v0.5.0 \\
  & Infrastructure: Tests & \texttt{tests/} & $\sim$15,000 & v0.8.0 \\
  & Infrastructure: CLI Tools & \texttt{tools/} & $\sim$5,100 & v0.4.0 \\
  & Infrastructure: Utilities & \texttt{utils/}, \texttt{scripts/} & $\sim$500 & v0.9.7--v0.9.9 \\
\bottomrule
\end{tabular}
\end{table}

\begin{verbatim}
  PAI Oncology Trial FL - End-to-End Architecture (v1.1.0)
  ========================================================

  +-- Multi-Organization Cooperation (Pillar 5) ----------+
  |  FedAvg | FedProx | SCAFFOLD | Secure Aggregation     |
  |  Site Enrollment | Data Harmonization | DP Budget      |
  +-------+------+------+------+------+------+------------+
          |      |      |      |      |
  +-------+--+ +-+------+-+ +--+------+--+ +-+----------+
  | Privacy  | | Regulatory | | Unification | | Standards |
  | (Pillar  | | (Pillar 2) | | (Pillar 3)  | | (Pillar  |
  |  1)      | |            | |             | |  4)       |
  | PHI Det. | | FDA Track  | | Isaac-MuJo  | | Model     |
  | De-ID    | | IRB Mgmt   | | Bridge      | | Registry  |
  | RBAC     | | GCP Check  | | Agent Unif. | | Benchmark |
  | Breach   | | eCTD Comp. | | Format Conv | | Convert   |
  +----+-----+ +-----+------+ +------+------+ +-----+----+
       |             |               |               |
  +----+-------------+---------------+---------------+----+
  |        Supporting Platform Modules                     |
  | Digital Twins | Clinical Analytics | Agentic AI | CLI  |
  | PK/PD | Survival | Risk | DSMB | MCP | ReAct | RAG   |
  +-------------------------------------------------------+
          |
  +-------+-----------------------------------------------+
  |              Tests & Quality Infrastructure            |
  | 82 test files | 15 directories | Security Audit       |
  | Triple AI Peer Review | CI (3.10/3.11/3.12)           |
  +-------------------------------------------------------+
\end{verbatim}

\subsection{Pillar 1: Privacy Infrastructure (v0.6.0)}
\label{sec:privacy}

The privacy framework implements comprehensive HIPAA Safe Harbor compliance across five modules totaling $\sim$5,400 LOC within the \texttt{privacy/} directory.

\subsubsection{PHI/PII Detection and Management}

The \texttt{phi\_detector.py} module (v0.6.0) implements all 18 HIPAA Safe Harbor identifiers per 45 CFR 164.514(b)(2)~\cite{hipaa_safeharbor}:

\begin{lstlisting}[caption={PHI detection categories from \texttt{privacy/phi-pii-management/phi\_detector.py} (v0.6.0)}]
class PHICategory(Enum):
    NAME = "name"
    GEOGRAPHIC = "geographic_subdivision"
    DATE = "dates"
    PHONE = "phone_number"
    FAX = "fax_number"
    EMAIL = "email_address"
    SSN = "social_security_number"
    MEDICAL_RECORD = "medical_record_number"
    HEALTH_PLAN = "health_plan_beneficiary"
    ACCOUNT = "account_number"
    LICENSE = "certificate_license"
    VEHICLE = "vehicle_identifier"
    DEVICE = "device_identifier"
    URL = "web_url"
    IP_ADDRESS = "ip_address"
    BIOMETRIC = "biometric_identifier"
    PHOTO = "full_face_photo"
    OTHER = "other_unique_identifier"
\end{lstlisting}

\subsubsection{De-identification Pipeline}

The de-identification module provides six transformation methods with HMAC-SHA256 pseudonymization using cryptographically random salts (\texttt{os.urandom(32)}), as hardened during peer review cycle 3 (v0.9.8/v0.9.9):

\begin{itemize}
  \item \textbf{REDACT:} Complete value removal
  \item \textbf{HASH:} HMAC-SHA256 pseudonymization (environment-driven keys post-v0.9.9)
  \item \textbf{GENERALIZE:} Geographic data to 3-digit ZIP codes
  \item \textbf{SUPPRESS:} Selective field suppression
  \item \textbf{PSEUDONYMIZE:} Consistent patient-level token generation
  \item \textbf{DATE\_SHIFT:} Consistent patient-level date offsets
\end{itemize}

\subsubsection{Access Control and Audit}

Role-based access control (RBAC) implements 21 CFR Part 11 compliant audit trails~\cite{cfr_part11} with five role tiers, time-limited access with fail-closed expiration (v0.6.0), and immutable audit log copies returned by \texttt{get\_audit\_log()} (v0.7.0 security fix).

\subsubsection{Breach Response and Data Use Agreements}

The breach response protocol implements four-factor risk assessment per 45 CFR 164.402 with 60-day notification tracking. The DUA generator supports five agreement types (research, clinical, vendor, multi-institutional, pediatric) with security tiers and configurable retention policies.

\subsection{Pillar 2: Regulatory Infrastructure (v0.6.0, v0.9.1)}
\label{sec:regulatory}

\subsubsection{Core Regulatory Framework (v0.6.0)}

The \texttt{regulatory/} directory implements five compliance domains:

\begin{table}[H]
\centering
\caption{Regulatory infrastructure modules and compliance alignment}
\label{tab:regulatory}
\small
\begin{tabular}{@{}lll@{}}
\toprule
\textbf{Module} & \textbf{Version} & \textbf{Standards} \\
\midrule
\texttt{fda\_submission\_tracker.py} & v0.6.0 & FDA 510(k), De Novo, PMA \\
\texttt{irb\_protocol\_manager.py} & v0.6.0 & ICH E6(R3), 21 CFR Part 11 \\
\texttt{gcp\_compliance\_checker.py} & v0.6.0 & ICH E6(R3) 13 principles \\
\texttt{regulatory\_tracker.py} & v0.6.0 & FDA, EMA, PMDA, TGA, HC \\
\texttt{HUMAN\_OVERSIGHT\_QMS.md} & v0.6.0 & CRF risk tiers, AE boundaries \\
\bottomrule
\end{tabular}
\end{table}

The FDA submission tracker supports five regulatory pathways---510(k), De Novo, PMA, Breakthrough, and Pre-Submission---referencing FDA AI/ML Device Guidance (January 2025)~\cite{fda2025aiml}, PCCP Guidance (August 2025)~\cite{fda2025pccp}, and QMSR (February 2026). The GCP compliance checker scores 13 ICH E6(R3) principles~\cite{ich2025e6r3} excluding \texttt{NOT\_ASSESSED} findings from the denominator (v0.7.0 logic fix).

\subsubsection{Regulatory Submissions Platform (v0.9.1)}

The \texttt{regulatory-submissions/} platform ($\sim$4,600 LOC across 6 modules) complements clinical analytics with downstream regulatory packaging:

\begin{enumerate}
  \item \textbf{Submission Orchestrator:} Lifecycle management for FDA submission workflows
  \item \textbf{eCTD Compiler:} Module structure generation per FDA Technical Conformance Guide with SHA-256 checksums
  \item \textbf{Compliance Validator:} Multi-regulation validation across 9 frameworks (HIPAA, GDPR, FDA 21 CFR 820, Part 11, ISO 14971, IEC 62304, ICH E6(R3)/E9(R1), MDR 2017/745)
  \item \textbf{Document Generator:} Automated CSR, SAP, PCCP, ISO 14971, IEC 62304 documents
  \item \textbf{Regulatory Intelligence:} 7-jurisdiction guidance tracking (FDA, EMA, PMDA, TGA, Health Canada, MHRA, NMPA)
  \item \textbf{Submission Analytics:} KPIs, trend analysis, FDA MDUFA benchmarking
\end{enumerate}

\subsection{Pillar 3: Cross-Framework Unification (v0.3.0)}
\label{sec:unification}

The \texttt{unification/} directory (v0.3.0) enables seamless interoperability across four simulation frameworks for surgical robotics:

\begin{verbatim}
  Cross-Framework Unification Architecture (v0.3.0)
  ==================================================

  +------------------+       +------------------+
  | NVIDIA Isaac Sim |<----->|     MuJoCo       |
  | (High-fidelity)  |       |  (Fast physics)  |
  +--------+---------+       +--------+---------+
           |    Bidirectional Bridge   |
           +----------+  +------------+
                      |  |
               +------+--+------+
               | Unified State   |
               | Representation  |
               +------+--+------+
                      |  |
           +----------+  +------------+
           |                          |
  +--------+---------+       +--------+---------+
  |     Gazebo       |       |    PyBullet      |
  |   (ROS 2)        |       | (Lightweight)    |
  +------------------+       +------------------+
\end{verbatim}

Key unification components include:
\begin{itemize}
  \item \textbf{Isaac-MuJoCo Bridge} (\texttt{simulation\_physics/}): Bidirectional state conversion with parameter mapping, enabling policy transfer between high-fidelity and fast physics engines
  \item \textbf{Unified Agent Interface} (\texttt{agentic\_generative\_ai/}): Tool format conversion for CrewAI, LangGraph, AutoGen, and custom backends
  \item \textbf{Cross-Platform Tools} (\texttt{cross\_platform\_tools/}): Framework detection (Isaac Lab, MuJoCo, Gazebo, PyBullet), model format conversion (URDF, MJCF, SDF, USD), and policy export with validation
  \item \textbf{Standards Protocols} (\texttt{standards\_protocols/}): Data format standards, communication protocols, and safety compliance
\end{itemize}

\subsection{Pillar 4: Standards \& Benchmarking (v0.3.0)}
\label{sec:standards}

The \texttt{q1-2026-standards/} directory implements three Q1 2026 objectives:

\begin{enumerate}
  \item \textbf{Objective 1---Model Conversion Pipeline:} Cross-framework model format conversion with validation, supporting URDF$\leftrightarrow$MJCF$\leftrightarrow$SDF$\leftrightarrow$USD bidirectional transforms
  \item \textbf{Objective 2---Model Registry:} Federated model registry with validation gates, version tracking, and compliance metadata
  \item \textbf{Objective 3---Cross-Platform Benchmarking:} Standardized benchmark runner enabling reproducible comparison across simulation frameworks
\end{enumerate}

The implementation guide includes a timeline (\texttt{timeline.md}) and compliance checklist (\texttt{compliance\_checklist.md}) aligned with IEC 80601-2-77~\cite{iec80601} and ISO 14971~\cite{iso14971}.

\subsection{Pillar 5: Multi-Organization Cooperation (v0.1.0--v1.0.1)}
\label{sec:cooperation}

The federated architecture enables multi-organization cooperation across the oncology trial ecosystem:

\begin{table}[H]
\centering
\caption{Multi-organization cooperation model with platform integration points}
\label{tab:cooperation}
\small
\begin{tabular}{@{}p{3.2cm}p{3.8cm}l@{}}
\toprule
\textbf{Organization} & \textbf{Integration Points} & \textbf{Version} \\
\midrule
Academic Medical Centers & Federated training, IRB, digital twins & v0.1.0+ \\
Community Hospitals & Site enrollment, data harmonization & v0.2.0+ \\
Pharmaceutical Cos. & Treatment simulation, regulatory & v0.6.0+ \\
Robotics Manufacturers & Simulation frameworks, surgical tasks & v0.3.0+ \\
Regulatory Bodies & Submission tracking, compliance & v0.9.1+ \\
Standards Organizations & IEC/ISO/ICH alignment & v0.3.0+ \\
\bottomrule
\end{tabular}
\end{table}

The federated learning core (\texttt{federated/}) implements three aggregation strategies with privacy guarantees:
\begin{itemize}
  \item \textbf{FedAvg} (v0.1.0)~\cite{mcmahan2017fedavg}: Standard federated averaging for IID data distributions
  \item \textbf{FedProx} (v0.2.0)~\cite{li2020fedprox}: Proximal term regularization for non-IID heterogeneous hospital data
  \item \textbf{SCAFFOLD} (v0.2.0)~\cite{karimireddy2020scaffold}: Control variate correction for client drift across sites
  \item \textbf{Differential Privacy} (v0.1.0)~\cite{dwork2014dp}: Gaussian mechanism with gradient clipping and privacy budget ($\varepsilon$, $\delta$) tracking
  \item \textbf{Secure Aggregation} (v0.1.0): Mask-based protocol ensuring no raw model weights are exchanged
\end{itemize}

\subsection{Workflow Demonstrations: Domain Examples (v0.5.0)}
\label{sec:examples}

Version 0.5.0 introduced 31 production-quality example scripts across five domains, demonstrating end-to-end workflow capabilities. These examples serve as reference implementations that clinicians, engineers, and regulators can use to evaluate the platform.

\subsubsection{Core Federated Learning Examples}

Five core examples (\texttt{examples/}, 4,924 LOC total) demonstrate progressive complexity:
\begin{enumerate}
  \item \texttt{01\_federated\_training\_workflow.py} (881 LOC): Multi-site FedAvg/FedProx pipeline with DP integration, secure aggregation, and convergence monitoring
  \item \texttt{02\_digital\_twin\_planning.py} (924 LOC): Patient digital twin construction with exponential, logistic, Gompertz, and von Bertalanffy tumor growth models
  \item \texttt{03\_cross\_framework\_validation.py} (865 LOC): Policy validation across Isaac Sim, MuJoCo, PyBullet, and Gazebo with kinematic drift analysis
  \item \texttt{04\_agentic\_clinical\_workflow.py} (917 LOC): Multi-agent clinical decision support with oncologist, radiologist, safety, and data agents
  \item \texttt{05\_outcome\_prediction\_pipeline.py} (982 LOC): Federated outcome prediction with survival analysis, biomarker integration, and fairness auditing
\end{enumerate}

\subsubsection{Agentic AI Examples---Production Architecture}

Six agentic AI examples (\texttt{agentic-ai/examples-agentic-ai/}, $\sim$6,700 LOC total) provide production-grade implementations:

\begin{table}[H]
\centering
\caption{Agentic AI example scripts with architecture patterns (v0.5.0)}
\label{tab:agentic}
\small
\begin{tabular}{@{}clp{5cm}@{}}
\toprule
\textbf{\#} & \textbf{Script} & \textbf{Architecture Pattern} \\
\midrule
01 & \texttt{mcp\_oncology\_server.py} & Model Context Protocol server with tool registration, patient lookup, treatment simulation, protocol compliance, robotic telemetry \\
02 & \texttt{react\_treatment\_planner.py} & ReAct reasoning loops with thought-action-observation cycles, 10+ step traces \\
03 & \texttt{realtime\_adaptive\_monitoring.py} & Streaming multi-modal data, adaptive thresholds, CTCAE v5.0 grading \\
04 & \texttt{autonomous\_simulation\_orchestrator.py} & Multi-framework job coordination with Pareto optimization \\
05 & \texttt{safety\_constrained\_executor.py} & IEC 80601-2-77 safety gates, human-in-the-loop, emergency stop \\
06 & \texttt{oncology\_rag\_compliance.py} & RAG with semantic chunking, cross-encoder reranking, citation tracking \\
\bottomrule
\end{tabular}
\end{table}

\textbf{MCP Oncology Server (\texttt{01\_mcp\_oncology\_server.py}).} This script implements a Model Context Protocol server~\cite{mcp2025} exposing five oncology domain tools for LLM agents:

\begin{lstlisting}[caption={MCP tool registration pattern from \texttt{01\_mcp\_oncology\_server.py} (v0.5.0)}]
# Tool definitions with JSON Schema input validation
TOOLS = [
    {"name": "patient_lookup",
     "description": "Retrieve de-identified patient record",
     "inputSchema": {"type": "object",
       "properties": {"patient_id": {"type": "string"}}}},
    {"name": "treatment_simulation",
     "description": "Run digital twin treatment simulation",
     "inputSchema": {"type": "object",
       "properties": {"patient_id": {"type": "string"},
                      "modality": {"type": "string"}}}},
    {"name": "protocol_compliance",
     "description": "Verify ICH E6(R3) compliance",
     "inputSchema": {"type": "object",
       "properties": {"protocol_id": {"type": "string"}}}},
]
\end{lstlisting}

The MCP server integrates with the broader platform by:
\begin{itemize}
  \item Querying the \texttt{privacy/} de-identification pipeline before returning patient records
  \item Invoking the \texttt{digital-twins/} patient modeling engine for treatment simulation
  \item Validating against the \texttt{regulatory/} GCP compliance checker for protocol verification
  \item Logging all tool invocations to 21 CFR Part 11 compliant audit trails with hash-chain integrity
  \item Ingesting robotic procedure telemetry from physical AI surgical systems
\end{itemize}

\textbf{Synergistic inter-script operation.} The six agentic AI scripts function both independently and synergistically:
\begin{itemize}
  \item Scripts 01 (MCP) + 05 (Safety) combine to provide safety-constrained tool execution for LLM agents operating robotic systems
  \item Scripts 02 (ReAct) + 06 (RAG) combine reasoning with regulatory evidence grounding
  \item Scripts 03 (Monitoring) + 04 (Orchestration) enable real-time adaptive simulation campaigns where monitoring alerts trigger simulation parameter adjustments
  \item All six scripts share the platform's audit logging infrastructure, differential privacy protections, and HIPAA-compliant data handling
\end{itemize}
